
\subsection{Explaining geometric error: uncertainty signals}
\label{sec:unc_method}
To explain \emph{where} a reconstruction disagrees with the stereo reference, we compare per-pixel signals that are computed from a trained model and its monocular inputs only, and score them against the per-frame scale-aligned, stereo-reference-relative error $e(\mathbf{p})=|s_t\hat{D}(\mathbf{p})-z^{\mathrm{S}}(\mathbf{p})|/z^{\mathrm{S}}(\mathbf{p})$ of Section~\ref{sec:evalproto}. All signal maps are computed and frozen \emph{before} any stereo-derived mask is applied; the stereo reference only selects the pixels on which they are scored and defines the target.
\begin{itemize}
\item \emph{Accumulated opacity} $\hat{w}$, which Endo-4DGS uses as its confidence map (Section~\ref{sec:objective}); uncertainty $u_{\mathrm{cov}}=1-\hat{w}$.
\item \emph{Ray depth spread.} With the compositing weights $a_i$, the normalised raw moments $m_r=\sum_i a_i z_i^r/\hat{w}$ for $r\in\{1,2\}$ give $\sigma_z=\sqrt{\max(m_2-m_1^2,0)}$ and $u_{\mathrm{ray}}=\sigma_z/\max(\hat{D},\epsilon)$ with $\epsilon=10^{-6}$; $m_1, m_2$ are obtained in one extra rasterisation pass that renders $(z_i, z_i^2)$ as colours on a black background.
\item \emph{Deformation magnitude} $u_{\mathrm{mot}}=\big(\sum_i a_i\|\Delta\bm{\mu}_i\|/\hat{w}\big)/\max(\hat{D},\epsilon)$ at the frame's time.
\item \emph{Prior disagreement} $u_{\mathrm{prior}}=|\hat{D}-\tilde{z}|/\max(\hat{D},\epsilon)$, where $\tilde{z}=c\,z^{\mathrm{prior}}$ is the monocular prior of the same frame (Depth-Anything or MoGe-2, whichever the model was built from) and $c=\mathrm{median}(\hat{D}/z^{\mathrm{prior}})$ over pixels where both are finite and positive.
\item \emph{Seed ensemble} $u_{\mathrm{ens}}$: the relative standard deviation of the rendered depth of three models trained with different seeds, each normalised by its own median over the image.
\end{itemize}
We also test a combination without fitted parameters, the mean of the per-frame percentile ranks of $u_{\mathrm{ray}}$, $u_{\mathrm{prior}}$ and $u_{\mathrm{ens}}$, which unlike learned predictors~\cite{gottwald2025primu} needs no held-out error.

\textbf{Metrics.} Signals are scored on $\Omega\cap\{z^{\mathrm{prior}}>0\}$. For each test frame and seed we sort pixels by decreasing uncertainty, remove the most uncertain fraction $f$ for 50 values $f\in\{0,0.02,\dots,0.98\}$ and record the mean AbsRel of the remaining pixels; the oracle curve sorts by the true error. AUSE is the trapezoidal area between a signal's curve and the oracle curve, and AURG the area between the random-order curve (the constant mean error) and the signal's curve~\cite{poggi2020uncertainty}; both are unnormalised AbsRel integrals. We summarise each signal by the share of the attainable ranking gain,
\begin{equation}
G=100\,\frac{\operatorname{mean}_{t,s}\mathrm{AURG}_{t,s}}{\operatorname{mean}_{t,s}\mathrm{AUSE}^{\mathrm{random}}_{t,s}},
\label{eq:gain}
\end{equation}
a ratio of means over frames $t$ and seeds $s$, where $\mathrm{AUSE}^{\mathrm{random}}$ is the area between the random and oracle curves; $G{=}100\,\%$ matches the oracle and $G{=}0$ a random order. The mean of per-frame ratios differs from $G$ by at most 5.2 percentage points over all signals, setups and references and leads to the same conclusions. We also report Spearman's $\rho$ between $u$ and $e$ (ties by average rank), and the mean error per uncertainty decile, a ranking diagnostic rather than a calibration of predicted error magnitudes.

\textbf{Strict reference.} SGBM depth is itself an estimate. We therefore repeat every evaluation with a strict reference that keeps only pixels whose left- and right-view disparities agree within 1\,px (the right-view disparity is computed on the mirrored pair), retaining 20--25\,\% of the pixels. This reduces sensitivity to some stereo mismatches, but mutually consistent wrong matches remain and the evaluated pixels change.

\textbf{Failure attribution.} We partition the scored pixels into four disjoint regions with fixed precedence: the blue ring-shaped device visible in P1 (OpenCV hue 80--130, saturation $>0.35$, value $>0.15$, opened with a 5$\times$5 kernel), then specular highlights (value $>0.9$ and saturation $<0.3$, dilated with a 5$\times$5 kernel), then depth edges (top 10\,\% of $\|\nabla z^{\mathrm{S}}\|/z^{\mathrm{S}}$ over valid pixels, with 3$\times$3 Sobel gradients after Gaussian smoothing with $\sigma{=}1.5$\,px; invalid stereo pixels are filled with the median depth before smoothing, which can create spurious edges at mask borders), and the remaining tissue. Edge detection uses the reference, so attribution is an analysis tool rather than a test-time signal. We report the mean AbsRel per region and each region's share of the total error, computed from per-frame error sums.
